\chapter{Introduction}

\begin{center}
    ``When faced with a challenge, get smarter.'' \\
    - Ed Catmull \cite{Book-CreativityInc}
\end{center}

\par The Information Age has been one of the most challenging and revealing, yet fruitful times in the history of the world. For several decades in the early- to mid-twentieth century, a computer was considered a calculating machine, which consumed one or more rooms. The rooms, full of cabinets which contained blinking lights, controls, and the logic, arithmetic, and memory capabilities that makes computing possible, was seen as a wonder by many. Within the remainder of the twentieth century, the computer became known as a small device that could sit on someone's desk, in someone's lap, or within a cabinet in a data center. Within the first two decades of the twenty first century, though, the computer became much more than that. We now see them in our hands, on our heads, on our wrists, and beyond. Interestingly, over this nearly one hundred year period of expansion of the digital frontier, the computer is still fundamentally the same as it was in the beginning. The main difference now is the computer is magnitudes more performant and compact.

\par With changes in the size and scale of the computer comes changes in the means in which people interface with the computer. Human-Computer Interaction (HCI) - a field dedicated to the research, development, and innovation of how people use computing devices - has been a great influence in how the form factor and presentation of the computer has changed over time.\cite{Web-IDFHCI} Early HCI implementations include computer wiring patterns (e.g. ENIAC), punch cards (e.g. UNIVAC, IBM 1401), and bit manipulation by means of physical switches (e.g. Altair 8800, IMSAI 8080). In time, the keyboard was introduced, mimicking the familiar typewriters of the day.

\par However, in the late 1960s, the modern computer interface would be recognized for the first time - the mouse. Douglas Engelbart introduced radically new technology that floored the computer science community in what writer Steven Levy called ``The Mother of All Demos''. \cite{Book-InsanelyGreat} Unfortunately, it wouldn't be until years later when the keyboard and mouse (and its modern variants) would finally become the standard in interface for the computer.

\par Given the advances with HCI over time, we began seeing more interesting and experimental means of computer interface. Extended Reality (XR) technology - a blend of prominent new methodologies for interacting with a virtual or enhanced world - is one such means. Within XR, we see Virtual Reality (VR), Augmented Reality (AR), and Mixed Reality (MR). VR gives us the ability to create amazing, and sometimes even surreal or outright physically impossible, experiences contained within the bounds of hardware and software. AR provides a context where virtual elements are imposed into the physical world. MR offers the hybrid solution - virtual elements (seen either by means of VR or AR), being blended seamlessly into physical space. \cite{Web-HPXR}

\par While XR is seen by the general public as a means of providing entertainment value (i.e. games, film, artistic experiences, etc.), it can also be used as a means of enhancing and expanding our understanding of the world around us. Some other examples of XR in the world are the presentation of news and weather (i.e. title cards for news stories next to a news anchor, projections of weather conditions symbols onto maps) and online shopping experiences (i.e. websites where you can test fit eyeglasses \cite{Web-GlassesUSATryOnline}, clothing, etc.).

\par Although, one application of XR is less exposed to the public for many reasons, mostly around intellectual property protection - manufacturing process training and supplementation. Given that these processes are, in most cases, proprietary, there are some general cases that are shared across industries. For instance, Microsoft's HoloLens \cite{Web-MicrosoftHololens} and Google's Glass \cite{Web-GoogleGlass} platforms offers MR experiences which enhance the user's immediate working space - whether that is in an office building or the floor of a factory.

\par In this thesis, we'll explore an application of XR for automobile operation. We'll proceed to implement a solution that satisfies the application's needs.
